\section{Functional Requirements}

The Origin AI system must fulfill the following functional requirements to provide a comprehensive pedagogical experience:

\begin{itemize}
    \item \textbf{Contextual Awareness:} The system must capture and utilize the student's current page state, including the active question, subject, and chapter.
    \item \textbf{In-Situ Interaction:} Students must be able to select text on the screen and receive an explanation directly within a floating widget.
    \item \textbf{Multimodal Input:} The system must support both text-based queries and voice-based (transcribed) inputs.
    \item \textbf{Tiered Resolution:} Queries should be resolved using the most cost-effective method available (Stored Solution $\rightarrow$ Subject KB $\rightarrow$ Vector Retrieval $\rightarrow$ LLM).
    \item \textbf{Pedagogical Guardrails:} The system must enforce academic integrity by blocking answers during tests and providing hints instead of solutions for unattempted practice questions.
    \item \textbf{Subject Disambiguation:} The system must accurately route queries to the appropriate subject knowledge base (Biology, Chemistry, Mathematics, Physics).
    \item \textbf{Persistent Memory:} The system must maintain a history of student interactions and weak/strong topics to personalize responses.
\end{itemize}

\section{Non-Functional Requirements}

The system must also adhere to several non-functional requirements to ensure scalability and reliability:

\begin{itemize}
    \item \textbf{Performance (Latency):} The median response time for locally resolved queries should be under 500ms. For LLM-backed queries, the response should begin streaming within 3 seconds.
    \item \textbf{Scalability:} The backend must scale horizontally to handle at least 100 concurrent AI mentorship sessions.
    \item \textbf{Cost Efficiency:} The system should achieve at least a 75\% reduction in operational costs compared to a direct LLM-only approach.
    \item \textbf{Security:} Internal API endpoints must be protected from external access and require valid service tokens.
    \item \textbf{Availability:} The system should maintain 99.9\% uptime for core mentorship features.
\end{itemize}

\section{Scope and Non-Goals}

\subsection{In Scope}
\begin{itemize}
    \item Fully functional AI mentor integrated into the Origin platform.
    \item Support for four core subjects: Biology, Chemistry, Mathematics, Physics.
    \item Implementation of a 4-step tiered fallback pipeline.
    \item Persistent vector storage for session memory.
    \item Deployment on GCP Cloud Run and Vercel.
\end{itemize}

\subsection{Out of Scope}
\begin{itemize}
    \item Support for Indian languages beyond Hinglish in the current phase.
    \item Offline-only operation (requires on-device models).
    \item Full automated grading of high-stakes subjective assessments.
    \item Real-time multi-student collaborative AI sessions.
\end{itemize}

\section{User Personas}

To design a truly student-centric system, we identified three primary user personas that Origin AI serves:

\begin{enumerate}
    \item \textbf{Persona 1: The High-Stakes Aspirant (JEE/NEET)}  
    \textit{Needs:} Speed, accuracy, and rigorous adherence to the curriculum.  
    \textit{Behavior:} Spends 8-10 hours daily on the platform, asks complex conceptual questions, and needs immediate clarification during late-night study sessions when human tutors are unavailable.
    \item \textbf{Persona 2: The Foundation Student (K-12)}  
    \textit{Needs:} Simplified explanations, visual aids, and encouraging feedback.  
    \textit{Behavior:} Struggles with basic concepts, uses voice mode more frequently, and requires Socratic guidance (hints) to build confidence.
    \item \textbf{Persona 3: The Self-Learner (Remote Area)}  
    \textit{Needs:} Low data consumption, offline-first assets, and affordable access.  
    \textit{Behavior:} Has limited internet connectivity, relies on local knowledge bases (HC Verma, etc.), and values the tiered architecture's cost-efficiency.
\end{enumerate}

\section{Detailed Use Cases}

\subsection{Use Case 1: In-Situ Doubt Resolution}
\textbf{Primary Actor:} Student \\
\textbf{Precondition:} Student is practicing questions in the OG Code arena. \\
\textbf{Flow:}
\begin{enumerate}
    \item Student encounters a difficult term in a Physics question.
    \item Student highlights the term using the mouse/touch.
    \item A floating "Ask Origin AI" chip appears near the selection.
    \item Student clicks the chip; the AI mentor opens with the highlighted text as context.
    \item AI identifies the context (Chapter: Kinematics) and provides a targeted explanation.
\end{enumerate}
\textbf{Postcondition:} Student understands the term without leaving the practice page.

\subsection{Use Case 2: Academic Integrity Enforcement}
\textbf{Primary Actor:} Student \\
\textbf{Precondition:} Student is taking a live Mock Test. \\
\textbf{Flow:}
\begin{enumerate}
    \item Student tries to ask Origin AI for the solution to Question 15.
    \item The system detects `PageContext.mode == "test"`.
    \item The AI responds with a polite refusal: "I cannot assist you during a live test. Focus on your attempt, and we can discuss this in the post-test analysis."
\end{enumerate}
\textbf{Postcondition:} Academic integrity is maintained.

\section{Hardware and Software Requirements}

\subsection{Development Environment}
\begin{itemize}
    \item \textbf{OS:} macOS / Linux (Ubuntu 22.04+)
    \item \textbf{Languages:} TypeScript (Frontend), Python 3.10+ (Backend)
    \item \textbf{Frameworks:} Next.js 14 (App Router), FastAPI
    \item \textbf{Database:} PostgreSQL with `pgvector` extension
    \item \textbf{AI Models:} Google Gemini 1.5 Pro/Flash, OpenAI GPT-4o
\end{itemize}

\subsection{Production Environment (Deployment)}
\begin{itemize}
    \item \textbf{Cloud Provider:} Google Cloud Platform (GCP)
    \item \textbf{Compute:} Cloud Run (Serverless)
    \item \textbf{Frontend Hosting:} Vercel
    \item \textbf{Static Assets:} Cloud Storage / Vercel Blob
\end{itemize}

